\chapter{致\texorpdfstring{\quad}{}谢}

时光荏苒，三年的硕士研究生生涯已临近尾声。回首这段充满挑战与收获的旅程，我的心中充满感激。在此，我谨以至诚之心，向所有在我求学道路上给予指导、帮助和关怀的单位和个人，致以最诚挚的谢意。

首先，我要衷心感谢裴海龙教授，您以严谨的治学态度和渊博的专业知识为我指明了科研的方向。在我飞行试验不顺利时，算法设计不理想时，您总是能够一针见血的指出问题，并且耐心地引导我解决问题。您宽厚待人的品格也深刻地感染着我，涵道的每一次炸机，我都陷入深深的自责，您总是微笑示意说“又发现一个问题，这次解决了下次会飞的更好”。您的言传身教让我受益匪浅，也让我在科研道路上走得更稳。

同时，我要感谢“自主系统与网络控制”教育部重点实验室为我提供先进的实验平台和开放的学术氛围。没有实验室的帮助，我就难以完成我的科研工作。感谢实验室的李慧良老师，在实验室的日常工作事务中为我提供无私的帮助，您的日夜操劳我会铭记于心。感谢实验室的程子欢师兄和蒙超恒师兄，你们亦师亦友，是你们带我踏入涵道组的大门，从飞机组装到代码调试再到论文撰写，每一个环节你们都是不可或缺的。感谢张子杰、苏伟浩、崔启权、陈沛和黄培烨对我论文中实验的帮助，使我的实验过程更加顺利。

感谢莫鸿强教授，您对我的论文提出了宝贵的建议。感谢我在论文中引用的参考文献的所有作者们，在阅读了你们的文献后我深受启发，才得以完成我的论文。感谢国家重点研发计划资助 [2023YFB4704900]、航空科学基金项目资助 [20220056060001]、国家自然科学基金重大科研仪器研制项目[61527810]、中国高校产学研创新基金新一代信息技术创新项目 [2022IT046]、中央高校基本科研业务费专项资金资助等项目提供充足的科研经费支持。

感谢实验室的黄霖杰、徐鼎、殷中、叶欢霆、程子啸、路文龙、罗秦薇、鲍紫怡、朱健锋、唐鸣、陈嘉鹏等同门们，与你们的朝夕相处让我在科研的道路上不再孤单。

最深沉的感谢献给我的家人，感谢你们的理解、支持和鼓励，是你们让我在追求理想的道路上不断前行，步伐愈加坚定。

千言万语难表心中感怀，唯有将这份感激铭记于心。感谢所有在我求学路上给予帮助的人们，你们的名字未能一一具列，但每一份善意都已被珍藏于心。再次向所有帮助过我的人致以诚挚的谢意！

谨以此文，致敬来时路，寄望新征程。

\begin{minipage}[t]{0.945\textwidth}%
	\begin{flushright}
		单喜龙\\
		% \today\\	% 自动时间
		2025年5月28日\\	%固定时间
		于华南理工大学
		\par\end{flushright}
\end{minipage}

